\documentclass{article}
\newcommand{\bandname}{Grades 4--5}
% Shared preamble for the three packets. Compile with lualatex.
\usepackage[letterpaper,left=0.5in,right=0.5in,top=0.9in,bottom=0.65in,
            headheight=22pt,headsep=0.22in,footskip=0.32in]{geometry}
\usepackage{fontspec}
\setmainfont{Andika}
\usepackage{xcolor}
\usepackage{tikz}
\usepackage{fancyhdr}
\usepackage{graphicx}

\definecolor{pgreen}{RGB}{92,176,74}
\definecolor{pblue}{RGB}{70,112,205}
\definecolor{pred}{RGB}{222,66,52}
\definecolor{ppurple}{RGB}{128,82,172}
\definecolor{pyellow}{RGB}{246,203,40}

\tikzset{
  grid/.style={line width=0.5pt, black!38},
  outline/.style={line width=1.7pt, black, line join=round, line cap=round},
  kgrid/.style={line width=0.5pt, black!30},
  koutline/.style={line width=2pt, black, line join=round, line cap=round},
  cgrid/.style={line width=0.35pt, black!32},
  coutline/.style={line width=1pt, black, line join=round},
  boardfill/.style={fill=white},
  hole/.style={fill=black!45},
  piece/.style={line width=1.1pt, black, line join=round},
  pfill/.style={fill=pblue!22},
  ribbonfill/.style={fill=pblue!62},
  iconline/.style={line width=0.8pt, black, line join=round},
  boardlabel/.style={font=\bfseries\Large, inner sep=1pt},
  ribbonletter/.style={font=\bfseries\large, text=white},
  cubeL/.style={fill=pblue!40},
  cubeR/.style={fill=pblue!78},
  cubeS/.style={fill=pblue!7},
}

\pagestyle{fancy}
\fancyhf{}
\lhead{Bellingham Math Circle}
\chead{\bandname}
\rhead{Name \rule[-0.15ex]{2.1in}{0.5pt}}
\cfoot{\thepage}
\renewcommand{\headrulewidth}{0.5pt}

\setlength{\parindent}{0pt}
\setlength{\parskip}{0.55em}
\raggedright
\hyphenpenalty=10000
\exhyphenpenalty=10000

\newcommand{\prob}[1]{\textbf{Problem #1:}\ }
\newcommand{\ansbox}[1][0.85in]{\tikz[baseline=(b.base)]{\node[draw,line width=0.9pt,rounded corners=3pt,minimum width=#1,minimum height=0.55in](b){\strut};}}
\newcommand{\fig}[1]{\input{fig/#1.tex}}

\AtBeginDocument{\fontsize{13}{18}\selectfont}
\newcommand{\flipped}{\tikz[baseline=-0.5ex]{\draw[line width=1.2pt,<->,>=stealth] (0,0) -- (0.7in,0);}}

\begin{document}
\fontsize{13}{18}\selectfont

% ---------------------------------------------------------------- page 1
\prob{1}Cover board A with blue pieces.
Two coverings are different if some blue piece is in a different place.
Draw each covering you find in one of the copies of board A on page~2.
How many different coverings does board A have?

\vspace{0.1in}
\begin{center}
\hspace{0.3in}\fig{g45_board_A}
\end{center}

\vspace{0.25in}
\prob{2}When three blue pieces fill a small hexagon, you can lift them out and put them back the other way, as in the picture.
This move is a flip.
In the space on page~2, draw a map of the coverings of board A: one dot for each covering, with a line between two dots when one flip changes one covering into the other.
Which two coverings of board A are the most flips apart?

\vspace{0.1in}
\begin{center}
\raisebox{-0.5\height}{\fig{g45_flip_0}}\hspace{0.3in}\flipped\hspace{0.3in}\raisebox{-0.5\height}{\fig{g45_flip_1}}
\end{center}

\newpage
% ---------------------------------------------------------------- page 2
\begin{center}
\newcommand{\ca}{\fig{g45_copy_A}}
\begin{tabular}{c@{\hspace{0.35in}}c@{\hspace{0.35in}}c@{\hspace{0.35in}}c}
\ca & \ca & \ca & \ca \\[0.25in]
\ca & \ca & \ca & \ca
\end{tabular}

\vspace{0.35in}
\tikz{\draw[line width=0.6pt, black!50, rounded corners=6pt] (0,0) rectangle (7.2in,4.6in);}
\end{center}

\newpage
% ---------------------------------------------------------------- page 3
\prob{3}In a covering of board A, look at the blue piece that touches the top edge of the board.
Its bottom edge is flat, and another blue piece hangs below that edge.
Keep going down until you reach the bottom edge of the board.
This chain of pieces is the ribbon of the covering.
Going down, each piece of the ribbon steps to the left (L) or to the right (R).
The picture shows a covering of a different board with its ribbon shaded.
Its ribbon is RLRR.

Which strings of L and R are ribbons of coverings of board A?
Can two different coverings of board A have the same ribbon?

\vspace{0.05in}
\begin{center}
\fig{g45_ribbon_example}
\end{center}

\vspace{0.2in}
\begin{minipage}[t]{4.9in}
\raggedright
\prob{4}How many coverings does board B on page~4 have?

\vspace{0.1in}
How many coverings does a hexagon have if its top and bottom sides are 1~small edge long and its other four sides are 4~small edges long, like the one drawn here?
\end{minipage}\hfill
\begin{minipage}[t]{1.9in}
\vspace{0pt}
\hfill\fig{g45_sketch_144}
\end{minipage}

\newpage
% ---------------------------------------------------------------- page 4
\begin{center}
\hspace{0.3in}\fig{g45_board_B}
\end{center}

\vspace{0.15in}
\prob{5}Cover board C with blue pieces.
How many different coverings does board C have?
Draw the coverings you find in the copies of board C on page~5.

\vspace{0.05in}
\begin{center}
\hspace{0.3in}\fig{g45_board_C}
\end{center}

\newpage
% ---------------------------------------------------------------- page 5
\vspace*{0.1in}
\begin{center}
\newcommand{\cc}{\fig{g45_copy_C}}
\begin{tabular}{c@{\hspace{0.1in}}c@{\hspace{0.1in}}c@{\hspace{0.1in}}c@{\hspace{0.1in}}c}
\cc & \cc & \cc & \cc & \cc \\[0.45in]
\cc & \cc & \cc & \cc & \cc \\[0.45in]
\cc & \cc & \cc & \cc & \cc \\[0.45in]
\cc & \cc & \cc & \cc & \cc \\[0.45in]
\cc & \cc & \cc & \cc & \cc
\end{tabular}
\end{center}

\newpage
% ---------------------------------------------------------------- page 6
\prob{6}Build covering X on board C.
Change it into covering Y using flips.
What is the fewest number of flips that can do this?

\begin{center}
\begin{tabular}{c@{\hspace{1in}}c}
\fig{g45_X} & \fig{g45_Y}\\[2pt]
\textbf{X} & \textbf{Y}
\end{tabular}
\end{center}

\prob{7}Which two coverings of board C are the most flips apart?
How many flips apart are they?

\vspace{0.1in}
\prob{8}Start from a covering of board A, make flips one at a time, and stop when you are back at the covering you started with.
Can the number of flips be odd?
Answer the same question for board C.

\vspace{0.1in}
\begin{minipage}[t]{4.6in}
\raggedright
\prob{9}The picture shows a covering of board C, turned a quarter turn, with each piece shaded by the direction it points.
It can be seen as cubes stacked in the corner of a box.
What does one flip do to the stack of cubes?
How many cubes does the box hold when it is full?
\end{minipage}\hfill
\begin{minipage}[t]{2.5in}
\vspace{0pt}
\hfill\fig{g45_cubes}
\end{minipage}

\end{document}
